\documentstyle[%


righttag%


,11pt%


,amssymb%


,russian%               remove in Eng.


,izvru%                 Set izven% in Eng.


]{amsart}





%      Math definitions





\newcommand{\lu}{\mathrm{l}}


\newcommand{\ru}{\mathrm{r}}


% NB: change in Eng.^^ instead of л & п


\newcommand{\col}{\operatorname{col}}


\newcommand{\vraisup}{\operatorname{vrai~sup}}


\newcommand{\tts}[1]{}


\numberwithin{equation}{section}





%\hoffset=-15mm                  For Eng.





\setcounter{page}{55}


\god{1998}


\nomer{3~(430)} %                Remove in Eng.


%\nomer{Vol.\,42, No.\,2}        Set in Eng.


% \str{\pageref{begin}--\pageref{end}}  Set in Eng.


\udk{517.958}


\author{\it С.Л.\,Тонконог, Л.Д.\,Эскин}


% NB    ^^ remove in Eng.


\title{О некоторых симметриях и инвариантных решениях\\


уравнения динамики неньютоновской жидкости. II}





\thanks{Работа выполнена при финансовой поддержке Российского фонда


фундаментальных исследований (проект \No\,97-01-00346).}


\date{}


% \pagestyle{myheadings}         Set in Eng.





\pagestyle{plain} %              Remove in Eng.


\begin{document}





% \thispagestyle{plain}          Set in Eng.





\maketitle


\label{begin}


Предлагаемая статья является второй частью работы авторов [1], в


которой были классифицированы уравнения, описывающие динамику свободной


поверхности неньютоновской жидкости со степенным реологическим законом


и допускающие группу симметрий. Здесь изучаются инвариантные решения


найденных там уравнений, сохранены все обозначения из [1], нумерация


разделов и формул продолжает принятую там нумерацию. Список цитированной


литературы для удобства читателя приведен в [1].


\addtocounter{section}{2}





\section{Инвариантные решения}


\addtocounter{section}{3}





6. Построим решение уравнения (3.8), инвариантное относительно группы с


оператором (3.9). Это решение представляется в виде


\begin{equation}


u=t^{-(l+s)}\chi_1(J),\quad \sigma=t^{-3s}\chi_2(J)


\end{equation}


где инвариант $J=t(x-(\ln t-1+\varepsilon^{-1})\beta^{-1})$. После


подстановки выражений (6.1) в уравнение


(3.8) для неизвестных функций $\chi_1$, $\chi_2$ получим системы


обыкновенных дифференциальных уравнений


$$


(\chi^2_1\chi^n_2)'-(1+\varepsilon)J\xi'_1+(l+s)\chi_1+


\psi(\chi_1,\chi_2)=0,\quad \chi_2=\chi_1\chi'_1.


$$


При $\varepsilon=-3l$, $\psi=0$ имеем замкнутое решение ($n$ нечетно)


\begin{equation}


u=


\begin{cases}


c_nt^{-(l+s)}(\xi^p_0-J^p)^{l-s}, &|J|\le\xi_0,\ \;


c_n=(l^{-1}(l+s)^{1/n})^{l-s};\\


\quad 0, & |J|>\xi_0.


\end{cases}


\end{equation}


Решение (6.2) описывает течение пленки неньютоновской жидкости с


нулевыми стоками в точках фронта, движущихся по логарифмическому закону


$$


x^л_f=\beta^{-1}\biggl(\ln t-\frac{5n+4}{3n+3}\biggr)


-\xi_0t^{-1},\quad


x^п_f=\beta^{-1}\biggl(\ln t-\frac{5n+4}{3n+3}\biggr)


+\xi_0t^{-1}.


$$


Таким образом, вся масса пленки при $t\to\infty$ медленно движется


вправо, причем


ширина пленки (расстояние между точками фронта, равное $2\xi_0t^{-1}$)


быстро уменьшается. Быстро спадает со временем и высота купола пленки


$$


u_{\max}=c_n\xi_0^lt^{-(l+s)},


$$


причем сам купол также движется вправо по закону


$$


x_k=\beta^{-1}\biggl(\ln t-\frac{5n+4}{3n+3}\biggr).


$$


Заметим, что в случае нулевого баланса массы свободное растекание


пленки (растекание с нулевыми точками в точках фронта) описывается


решением (2.5) уравнения (2.1). В этом случае купол неподвижен


($x_k=0$),


его высота спадает медленнее ($u_{\max}\sim t^{-r}$), точки фронта


движутся к $\infty$ по


степенному закону. Все указанные особенности течения, описываемого


решением (6.2) уравнения (3.8) с ненулевым балансом массы


\begin{equation}


g=\varepsilon t^{-1}(\beta^{-1}(\ln t-1)-x)\sigma u^{-1}=


-\varepsilon(Jt^{-2}+(\varepsilon\beta t)^{-1})


\sigma u^{-1},\quad \varepsilon=-3l<0,


\end{equation}


можно объяснить следующим образом (полагая для определенности


$\beta>0$). Из


(6.2) имеем $\sigma u^{-1}>0$ при $J<0$, т.е. при


$x^п_f<x<x_k$, и $\sigma u^{-1}<0$ при $J>0$, т.е.


при $x^п_f>x>x_k$. Отсюда и


из (6.3) убеждаемся, что баланс массы $g$ на решении (6.3) отрицателен


при $x<x_k$, положителен при $x_k<x<\beta^{-1}(\ln t-1)$ и снова


отрицателен при $\beta^{-1}(\ln t-1)<x<x^п_f$, причем


при больших $t$ последний случай невозможен. Отрицательный баланс $g$


равносилен стоку на поверхности, а положительный --- источнику. Наличие


стока на поверхности слева от $x_k$ и приводит к тому, что пленка


движется


вправо. На временах $t<-\varepsilon\beta\xi_0$ сток превалирует над


источником, что приводит к быстрому опусканию купола.





При $\varepsilon=-1$ вместо (6.2) нетрудно получить решение в


квадратурах, но с


ненулевым стоком на фронте. Такие решения здесь не рассматриваются.





Решение уравнения (3.10), инвариантное относительно группы с


оператором (3.11), следует искать в виде


\begin{equation}


u=\chi_1(J),\quad \sigma=t^{-2m}\chi_2(J),


\end{equation}


где $J=xt^{-2m}+(\varepsilon m)^{-1}t^{-m}$


--- инвариант группы (3.11).


Подстановка в уравнение (3.10)


приводит к обыкновенным дифференциальным уравнениям для неизвестных


функций $\chi_1$, $\chi_2$


\begin{equation}


(\chi^2_1\chi^2_2)'+\varepsilon(m^2J^2\chi'_1+\psi


(\chi_1,\chi_2)=0,\quad \chi_2=\chi_1\chi_1'.


\end{equation}


В случае $\psi=0$, или, более общо, $\psi=\beta\chi^\mu_1\chi^\nu_2$,


где


$$


\mu+\frac{3}{n+2}\nu=\frac{3(n+1)}{n+2},


$$


уравнение (6.5) допускает группу растяжений с оператором


$$


J\frac{\partial}{\partial J}+(l+s)\chi_1


\frac{\partial}{\partial\chi_1},


$$


следовательно, допускает понижение порядка, т.е. сводится к


динамической системе на плоскости, которая может быть исследована


методами качественной теории обыкновенных дифференциальных уравнений.


\medskip


\addtocounter{section}{1}


\addtocounter{equation}{-5}





7. Рассмотрим


инвариантные решения уравнений п.\,4. В


случае уравнения (4.11)  решение,  инвариантное  относительно


группы с оператором (4.12), имеет вид


\begin{equation}


u=t^\alpha\chi_1(J),\quad \sigma=t^{-\lambda\alpha}\chi_2(J),


\end{equation}


где инвариант


$J=xt^{-\alpha(2+\lambda)}-\varepsilon t^{k-1}/(k-1)$, а функции


$\chi_1$, $\chi_2$ удовлетворяют уравнению


\begin{equation}


(\chi^2_1\chi^n_2)'+\alpha(2+\lambda)J\chi'_1-


\alpha\chi_1+\varepsilon\psi(\chi_1,\chi_2)=0,\quad


\chi_2=\chi_1\chi'_1.


\end{equation}


При $\psi=0$, $\lambda=-3$ оно имеет замкнутое решение, которое


получается из решения


(2.5) уравнения (2.1),  если в нем заменить автомодельную переменную


$\xi$


инвариантом  $J$. Это решение описывает


растекание пленки с подвижным куполом и при $\varepsilon=0$


превращается в


решение (2.5).





Наиболее интересно среди инвариантных уравнений п.\,4


уравнение (4.13). Его решение, инвариантное относительно группы


с оператором (4.14), получается в виде


\begin{equation}


u=t^\alpha\exp(\varepsilon l\mu(t))\chi_1(J),\quad


\sigma=t^{-\alpha\lambda}\exp(\varepsilon s\mu(t))\chi_2(J),


\end{equation}


где $\mu(t)=(k-1)^{-1}t^{k-1}$,


$J=xt^{-\alpha(2+\lambda)}\exp(-\varepsilon\mu(t))$ --- инвариант


группы с оператором (4.14), а функции $\chi_1$, $\chi_2$ удовлетворяют


уравнению


\begin{equation}


(\chi^2_1\chi^n_2)'+\alpha(2+\lambda)J\chi'_1-\alpha\chi_1


+\varepsilon J^l\psi(\chi_1J^{-l},\chi_2J^{-s})=0,\quad


\chi_2=\chi_1\chi'_1.


\end{equation}


В случае $\psi=0$, $\lambda=-3$  уравнение


(7.4)  имеет  замкнутое решение (аналог решения (2.5) уравнения


(2.1))


\begin{gather}


u=


\begin{cases}


c_n t^{-r}\exp(\varepsilon l\mu(t))(\xi^p_0-J^p)^{l-s}, &


|J|\le \xi_0;\\


\quad 0,& |J|\ge \xi_0,


\end{cases}


\\


J=xt^{-r}\exp(-\varepsilon\mu(t)).\notag


\end{gather}


Постоянная $c_n$  определена соотношением (2.5). Оказывается, что


решение (7.5) уравнения  (4.13)  описывает  течение  пленки или


ледника, существенно отличающееся по своим физическим свойствам


от течения, описываемого решением (2.5) уравнения с нулевым балансом


массы (2.1). Действительно, пусть $\varepsilon<0$. Из (7.5) имеем


\begin{align}


&x^л_f=-\xi_0t^r\exp(\varepsilon\mu(t)),\quad


x^п_f=\xi_0 t^r\exp(\varepsilon\mu(t)),\\


&u_k=u_{\max}=u(0,t)=c_n\xi^l_0t^{-r}


\exp(\varepsilon l\mu(t)).


\end{align}


Рассмотрим три случая.





1) $k>1$. В этом случае $x^п_f(t)$ растет ($x^л_f(t)$ убывает) на


интервале  $0<t<t_0=(-r/\varepsilon)^{1/(k-1)}$ от нуля до


максимального


значения, равного $\xi_0t^r_0\exp(-r(k-1)^{-1})$, затем с ростом $t$


убывает от своего максимального значения


($x^л_f$  растет от своего


минимального значения) до нуля. При этом высота купола  $u_k$  убывает


$\lim u_k=0$ ($t\to\infty$).


Таким образом, на временном интервале $0<t<t_0$


имеем растекание пленки (наступление ледника) и


последующее ее схлопывание (отступление ледника).





2) $k<1$. В этом случае в силу (7.6)


$x^п_f\to -\infty$, $x^л_f\to\infty$ при $t\to\infty$,


причем из (7.7) получаем $u_k\to 0$ при $t\to\infty$. Решение (7.5) с


ростом  $t$ приближается к решению (2.5) и также описывает растекание


пленки.





3) $k=1$. В этом случае вместо формул (7.3) получаем


$$


u=t^{\alpha+\varepsilon l}\chi_1(J),\quad


\sigma=t^{-\alpha\lambda+\varepsilon s}\chi_2(J),


$$


а функции $\chi_1$, $\chi_2$


снова удовлетворяют уравнению (7.4). В случае $\psi=0$, $\lambda=-3$


имеем замкнутое решение


\begin{equation}


u=


\begin{cases}


c_nt^{-r+\varepsilon l}(\xi^p_0-J^p)^{l-s}, &


|J|\le\xi_0;\\


\quad 0, & |J|\ge \xi_0.


\end{cases}


\end{equation}


Формулы (7.6), (7.7)


заменяются соотношениями


\begin{equation}


x^л_f=-\xi_0t^{r+\varepsilon},\quad


x^п_f=\xi_0t^{r+\varepsilon},\quad


u_k=c_n\xi^l_0 t^{-r+\varepsilon l}.


\end{equation}


Из соотношений (7. 9)


следует,  что  решение  (7.8)    уравнения


(4.13) описывает растекание пленки при $\varepsilon>-r$ и схлопывание


при


$\varepsilon<-r$, при $\varepsilon=-r$


точки фронта неподвижны, но купол  опускается.





Отмеченные особенности  течения  снова  можно   объяснить,


рассматривая баланс массы


$g=\varepsilon t^{k-2}(lu-x\sigma u^{-1})$, $\varepsilon<0$,


на замкнутых


решениях (7.5), (7.8). Баланс $g$  оказывается отрицательным,


т.е. играет роль стока. При $k>1$ и $k=1$, $\varepsilon<-r$ сток


достаточно


силен и приводит  к  схлопыванию  пленки, а в случае растекания


наличие слабого стока с поверхности приводит к  более быстрому,


чем в отсутствие стока, оседанию купола и  замедлению  движения


точек фронта.





Мы привели пример решения уравнения (4.13) в  элементарных


функциях --- решение (7.5).  Это решение описывает течение пленки


с нулевыми стоками в точках фронта. Нетрудно привести аналогичные


примеры  и  для  ряда  других  инвариантных  уравнений вида


(1.1), полученных в настоящей работе. Мы здесь на этих примерах


останавливаться не  будем,  а укажем возможность построения для


этих уравнений решений в  квадратурах  с  ненулевым  стоком  на


фронте. А именно, если снова положить $\psi=0$, но $\lambda=-2$


(следовательно,


$\alpha=-s$), то в уравнении (7.4) исчезнут слагаемые, содержащие


независимую


переменную $J$. Полученное после этого уравнение  допускает понижение


порядка и интегрируется в  квадратурах. Его неотрицательное монотонно


возрастающее решение $\chi^{(1)}_1$ определяется из уравнения ($n$


нечетно)


$$


J-\xi_0=\int^{\chi_1}_0


\frac{v^{(n+2)/n}dv}{(\xi_1-lrv^{1/nr})^{1/(n+1)}}=


M(\chi_1,\xi_1),\quad \xi_0\le J\le J_1=\xi_0+


M((\xi_1/(lr))^{nr},\xi_1),


$$


$\xi_0$, $\xi_1>0$ --- константы. Неотрицательное монотонно убывающее


решение $\chi^{(2)}_1$ определяется из уравнения


$$


\eta_0-J=\int^{\chi_1}_0


\frac{v^{(n+2)/n}dv}{(\eta_1-lrv^{1/nr})^{1/(n+1)}}=


M(\chi_1,\eta_1),\quad


\eta_0-M(\eta_1/(lr))^{nr},\eta_1)=J_2\le J\le \eta_0,


$$


$\eta_0$, $\eta_1>0$ --- константы. Имеем


\begin{align*}


&\chi^{(1)}_1(J_1)=(\xi_1/(lr))^{nr},\quad


\chi^{(2)}_1(J_2)=(\eta_1/(lr))^{nr},\\


&(\chi^{(1)}_1(J))'=(\xi_1-lr(\chi^{(1)}_1(J))^{1/(nr)})^{1/(n+1)}


(\chi^{(1)}_1(J))^{-(n+2)/n}, \\


&(\chi^{(2)}_1(J))'=-(\eta_1-lr(\chi^{(2)}_1(J)^{1/(nr)}


)^{1/(n+1)}(\chi^{(2)}_1(J))^{-(n+2)/n},


\end{align*}


откуда получаем


$$


\lim_{J\to J_1-0}(\chi^{(1)}_1(J))'=0,\quad


\lim_{J\to J_2+0}(\chi^{(2)}_1(J))'=0.


$$


Потребовав выполнения равенств $J_1=J_2$ и


$\chi^{(1)}_1(J_1)=\chi^{(2)}_1(J_1)$, найдем


$\eta_1=\xi_1$, $\eta_0=2J_1-\xi_0$.


Теперь на  интервале   $(\xi_0,\eta_0)$  без труда


строится неотрицательное непрерывное вместе с первой производной


решение $\chi_1$  уравнения (7.4) с $\lambda=-2$,  $\psi=0$  в виде


$$


\chi_1(J)=


\begin{cases}


\chi^{(1)}_1(J), &\xi_0\le J\le J_1;\\


\chi^{(2)}_1(J), & J_1\le J\le \eta_0.


\end{cases}


$$


Решение $\chi_1$  полностью определяется выбором постоянных $\xi_0$,


$\xi_1$.


Оно обращается в нуль при $J=\xi_0$  и $J=\eta0$, а его производная


равна нулю  в точке $J=J_1$. Имея решение $\chi_1$, с помощью


соотношения (7.3)


получаем решение уравнения (4.13) (с $\alpha=-s$)


с двумя точками фронта.


Нетрудно заметить,  что для  этого   решения


сток


в точке фронта $x^п_f=(2J_1-\xi_0)\exp(\varepsilon\mu(t))$ отличен от


нуля и


равен


$$


q=-t^{-2l}\exp(\varepsilon s\mu(t)/r)\xi_1^{n/(n+1)},


$$


а в точке фронта $x^л_f=\xi_0\exp(\varepsilon\mu(t))$


получаем отличный от


нуля источник


$$


q=t^{-2l}\exp(\varepsilon s \mu(t)/r)\xi_1^{n/(n+1)},


$$


производная $u_x$ при $x=x_k=J_1\exp(\varepsilon\mu(t))$ (координата


купола)


обращается в нуль.





Аналогичные решения  нетрудно построить и для других инвариантных


уравнений.





Автомодельное решение уравнения (4.16), инвариантное относительно


группы растяжений $X^0_\lambda$, ищем в виде (2.3), но для неизвестных


функций


$\chi_1$, $\chi_2$


из (4.16) вместо уравнения (2.4) получаем уравнение


$$


(\chi^2_1\chi^n_2)'+\alpha(2+\lambda)I_1\chi'_1-\alpha\chi_1+


\varepsilon I_1^{-(n+1)\lambda/(2+\lambda)}\psi(\chi_1


I_1^{-1/(2+\lambda)},\chi_2I_1^{2/(2+\lambda)},\quad


\chi_2=\chi_1'\chi_1.


$$


Уравнение (4.16) наиболее интересно в случае $\lambda=-1$  и функции


$\psi(p,q)$, зависящей лишь от отношения своих аргументов, т.е.


в случае уравнения


\begin{equation}


u_t=(u^2\sigma^n)_x+\varepsilon x^{n+1}\varphi(\sigma/u),\quad


\sigma=uu_x.


\end{equation}


Уравнения вида  (7.10)  имитируют  динамику поверхности ледника


в случае,  когда осадки на его поверхность  переносятся  ветром


постоянного направления (причем выбор функции $\varphi$ определяется


именно направлением ветра и мощностью  осадков,  предполагаемой


постоянной).





Перейдем к уравнению (4.21). Его решение, инвариантное относительно


группы с оператором (4.22), записывается в виде


$$


u=\chi_1(x)(t-\varepsilon\sqrt{\beta t^2+\beta_2})^{-s},\quad


\sigma=\chi_2(x)(t-\varepsilon\sqrt{\beta t^2+\beta_2})^{-2s},


$$


причем функции $\chi_1$, $\chi_2$  удовлетворяют уравнению


$$


(\chi^2_1\chi^n_2)'+s(1-\beta\varepsilon^2)\chi_1+


\varepsilon\chi_1^{2(n+1)}\psi(x,\chi_2/\chi^2_1),\quad


\chi_2=\chi_1\chi'_1.


$$





Рассмотрим теперь уравнение (4.26) при $c_3=0$ и $\theta_2=0$.


Уравнение (4.25)  удовлетворяется, а функции $\varphi$ и $\theta_1$


удовлетворяют


системе (4.24). Инвариантное относительно группы с оператором


$X_\lambda$(2.12)  решение уравнения (4.26) ищем в виде


\begin{equation}


u=\chi_1(J)\exp\biggl(\alpha\int\frac{1+\varepsilon\zeta_1}


{t+\varepsilon\alpha\varphi}dt\biggr),\quad


\sigma=\chi_2(J)\exp\biggl(\alpha\int


\frac{-\lambda+\varepsilon\zeta_2}{t+\varepsilon\alpha\varphi}dt


\biggr),


\end{equation}


где инвариант


$$


J=x\exp\biggl(-\alpha\int


\frac{2+\lambda+\varepsilon\theta_1}{t+\varepsilon\alpha\varphi}dt


\biggr).


$$


Подстановка соотношений (7.11) в уравнение (4.26) ($c_3=\theta_2=0$)


после достаточно  громоздких  преобразований, основанных на


использовании соотношений (4.27), приводит к уравнению для $\chi_1$,


$\chi_2$


$$


(\chi^2_1\chi^n_2)'+(\alpha(2+\lambda)+\varepsilon


(c_2+\varepsilon\alpha d_2))J\chi'_1-\alpha


[1+\varepsilon(l(\alpha^{-1}c_2+\varepsilon d_2)


-s(\varkappa+\varepsilon\alpha d_1))]\chi_1=0,\quad


\chi_2=\chi_1\chi'_1.


$$





В случае  уравнения  (4.29) решение,  инвариантное относительно группы


с оператором (4.30), получаем в виде


$$


u=t^{(\alpha+\varepsilon(l-s\mu))/(1+\varepsilon\mu)}\chi_1(J),\quad


\sigma=t^{(\varepsilon s(1-2\mu)-\alpha\lambda)/(1+\varepsilon\mu)}


\chi_2(J),\quad


J=xt^{-(\alpha(2+\lambda)+\varepsilon)/(1+\varepsilon\mu)}.


$$


Функции $\chi_1$, $\chi_2$ удовлетворяют уравнению


$$


(\chi^2_1\chi^n_2)'+


\biggl(\frac{\alpha(2+\lambda)+\varepsilon}{1+\varepsilon\mu}


+\varepsilon c_2\biggr)J\chi'_1-


\biggl(\frac{\alpha+\varepsilon(l-s\mu)}{1+\varepsilon\mu}


-\varepsilon c_1\biggr)\chi_1+ J^l\psi(\chi_1J^{-l},


\chi_2J^{-s})=0,\quad \chi_2=\chi_1\chi'_1.


$$





Приведем и пример решения уравнения (4.29),  инвариантного


относительно двупараметрической группы с операторами   (4.30) и


$X^0_\lambda$


$$


u=m_1 t^{-s}x^l,\quad \sigma=m_2t^{-2s}x^s,


$$


причем коэффициенты $m_1$, $m_2$ должны удовлетворять алгебраической


системе уравнений


$$


sr^{-1}m^2_1m^n_2+(s+\varepsilon c_1)m_1+


\varepsilon c_2 m_2 m^{-1}_1+\varepsilon\psi(m_1,m_2)=0,\quad


m_2=lm^2_1.


$$





Далее, с помощью (4.32) находим,  что инвариантное решение


уравнения (4.31) следует искать в виде


$$


u=t^\alpha\chi_1(J),\quad


\sigma=t^{-\lambda\alpha}\chi_2(J),\quad


J=xt^{-\alpha(2+\lambda)}-


\frac{\varepsilon\delta}{1+\varepsilon\beta}\ln t.


$$


Для функций $\chi_1$, $\chi_2$ получаем уравнение


$$


(\chi^2_1\chi^n_2)'+\alpha(2+\lambda)(1+\varepsilon\beta)


J\chi'_1+\varepsilon\biggl(c_3+


\frac{\delta}{1+\varepsilon\beta}\biggr)\chi'_1-


(\alpha-\varepsilon c_1)\chi_1+\varepsilon\psi(\chi_1,\chi_2)=0,\quad


\chi_2=\chi_1\chi'_1.


$$





В заключение укажем автомодельное решение уравнения


(4.33), оно задается соотношениями (2.3), но для функций


$\chi_1$, $\chi_2$


получаем уравнение


$$


(\chi^2_1\chi^n_2)'+[(\alpha(2+\lambda)+\varepsilon c_2)\xi


+\varepsilon c_3]\chi'_1-(\alpha-\varepsilon c_1)\chi_1+


\varepsilon\psi(\xi,\chi_1,\chi_2)=0,\quad


\chi_2=\chi_1\chi'_1.


$$


\medskip


\addtocounter{section}{1}


\addtocounter{equation}{-11}





8. Перейдем к построению инвариантных решений


уравнений п.\,5. Для уравнения (5.4),


инвариантного  относительно группы с оператором (5.5),


решение следует искать в  виде


\begin{equation}


u=\chi_1(J)\exp\Bigl(\varepsilon l\int\theta_1dt\Bigr),\quad


\sigma=\chi_2(J)\exp\Bigl(\varepsilon s\int\theta_1dt\Bigr),


\end{equation}


где


$$


J=\biggl(I_1-\lambda\delta^{-1}\Bigl(\int\theta_1dt+


\varepsilon^{-1}\Bigr) \biggr)\exp\Bigl(-\varepsilon\int


\theta_1dt\Bigr)+\lambda(\varepsilon\delta)^{-1}


$$


--- инвариант группы (5.5), равный $I_1=x+\lambda t$ при


$\varepsilon=0$.


После подстановки (8.1) в уравнение (5.4)


получим  для  определения неизвестных функций $\chi_1$, $\chi_2$


уравнение


\begin{equation}


(\chi^2_1\chi^n_2)'+(\varepsilon\delta J-\lambda)\chi'_1+


\varepsilon\chi_1\psi(\chi_1\chi^{-n-1}_2)=0,\quad


\chi_2=\chi_1\chi'_1.


\end{equation}


При $\varepsilon=0$ уравнение (8.2) переходит в уравнение для решений


типа бегущей


волны невозмущенного уравнения (2.1), имеющих вид (2.6). Отметим, что


уравнение (5.4)


содержит   произвольную функцию времени $\theta_1$,


что  позволяет получать решения (8.1) с достаточно разнообразными


свойствами и поведением на больших временах. В случае $\psi=\delta$


нетрудно


построить замкнутое решение для уравнения (8.2), а затем получить и


замкнутое решение уравнение (5.4)


$$


u=


\begin{cases}


c_n(\xi^p_0-(J-\lambda/(\varepsilon\delta))^p)^{l-s}


\exp\Bigl(\varepsilon l{\displaystyle \int}\theta_1dt\Bigr),&


|J-\lambda/(\varepsilon\delta)|\le\xi_0,\ \;


c_n=(\varepsilon\delta(l-s)^{-n})^s;\\


\quad 0, & |J-\lambda/(\varepsilon\delta)||>\xi_0.


\end{cases}


$$


Это решение обращается в нуль при $\varepsilon=0$.





Решение уравнения (5.6), инвариантное относительно группы с оператором


(5.7), имеет вид


\begin{equation}


u=\chi_1(J)\exp\biggl(\varepsilon l


\frac{\theta_0t}{1+\varepsilon\varphi_0}\biggr),\quad


\sigma=\chi_2(J)\exp\biggl(\varepsilon s


\frac{\theta_0t}{1+\varepsilon\varphi_0}\biggr),


\end{equation}


где функции $\chi_1$, $\chi_2$ удовлетворяют уравнению


\begin{multline}


(\chi^2_1\chi^n_2)'+\varepsilon\biggl(\delta-


\frac{\varepsilon\varphi_0\theta_0}{1+\varepsilon\varphi_0}\biggr)


\biggl(J-\frac{\lambda}{\varepsilon\delta}+


\frac{\lambda\varphi_0(\delta-\theta_0)}{\delta^2}\biggr)


\chi'_1+ \\


+\varepsilon\chi_1\biggl(\psi(\chi_1\chi^{-n-1}_2)


-\frac{l\theta_0}{1+\varepsilon\varphi_0}\biggr)=0,\quad


\chi_2=\chi_1\chi'_1,


\end{multline}


а


\begin{multline*}


J=\biggl[I_1+\lambda\delta^{-1}(\varphi_0-\varphi_0\theta_0


\delta^{-1}-\varepsilon^{-1}-\theta_0t)-


\frac{\varepsilon c\varphi_0}


{\delta+\varepsilon\varphi_0(\delta-\theta_0)}\exp


(\delta t/\varphi_0)\biggr]\times \\


\times\exp\biggl(-\frac{\varepsilon\theta_0t}{1+\varepsilon\varphi_0}\biggr)


+\lambda\delta^{-1}(\varepsilon^{-1}+\varphi_0\theta_0


\delta^{-1}-\varphi_0).


\end{multline*}


Инвариант $J$ снова выбран так, чтобы иметь $J=I_1$ при


$\varepsilon=0$.


Отметим, что в случае


$$


\psi=0,\quad \theta_0=


\frac{\delta(1+\varepsilon\varphi_0)}{\varepsilon\varphi_0-l}


$$


уравнение (8.4), а вместе с ним и (5.6), имеют замкнутое решение, на


деталях его построения не останавливаемся.





Перейдем к уравнению (5.8), инвариантному относительно группы с


оператором (5.9). Его инвариантное решение представляется в виде


\begin{equation}


u=(1+\varepsilon\nu t)^{-\mu/\nu}\chi_1(J),\quad


\sigma=(1+\varepsilon\nu t)^{-(2\mu+1)\nu^{-1}}\chi_2(J),


\end{equation}


где


$$


J=(1+\varepsilon\nu t)^{-1/\nu}(x-\lambda\varepsilon^{-1}-


\varepsilon b\nu t^{1/\nu})+\lambda\varepsilon^{-1}


$$


--- инвариант группы, удовлетворяющей условию


$J|_{\varepsilon=0}=I_1$, а функции $\chi_1$, $\chi_2$


удовлетворяют уравнению


\begin{equation}


(\chi^2_1\chi^n_2)'+\varepsilon((J-\lambda\varepsilon^{-1})\chi'_1+


\mu\chi_1+\psi(\chi_1\chi_2^{-\mu(2\mu+1)^{-1}})=0,\quad


\chi_2=\chi_1\chi'_1.


\end{equation}


В случае $\mu=1$, $\psi=0$ уравнение (5.8) снова решается в замкнутой


форме.





В заключение укажем инвариантное решение уравнения (5.10), полагая с


целью упрощения выкладок $\lambda=0$, $c=b(2n+2)^{-1}$ и заменив


произвольную постоянную $\mu$ на $-\mu$. Получаем уравнение


\begin{equation}


u_t=(u^2\sigma^n)_x-\varepsilon\biggl[\,


\frac{b}{2(n+1)}t^{-1/2}x+dt^{-1/2}\exp(2t^{1/2}b^{-1})


\sigma u^{-1}+\mu\sigma^{n+1}\biggr],\quad \sigma=uu_x,


\end{equation}


инвариантное относительно группы с оператором


$$


X=(1+\varepsilon bt^{1/2})\partial_t+\varepsilon\biggl[ x


\biggl(1+\frac{b}{2(n+1)}t^{-1/2}\biggr)+dt^{-1/2}\exp


(2t^{1/2}b^{-1})\biggr]\partial_x+


\varepsilon lu\partial_u+\varepsilon


\biggl(s-\frac{b}{2(n+1)}t^{-1/2}\biggr)\sigma\partial_\sigma.


$$


Инвариантное решение уравнения (8.7) имеет вид


\begin{align*}


u&=\chi_1(J)(1+\varepsilon bt^{1/2})^{-2l/(\varepsilon b^2)}


\exp(2lt^{1/2}b^{-1}),\\


\sigma&=\chi_2(J)(1+\varepsilon


bt^{1/2})^{-(\frac{1}{n+1}+\frac{2s}{\varepsilon b^2})}


\exp(2st^{1/2}b^{-1}),


\end{align*}


где


$$


J=(1+\varepsilon bt^{1/2})^{-\frac{1}{n+1}+


\frac{2}{\varepsilon b^2}}\biggl(x\exp(-2t^{1/2}b^{-1})-


\frac{2d}{b(-\frac{1}{n+1}+\frac{2}{\varepsilon b^2})}\biggr),


$$


а функции $\chi_1$, $\chi_2$ удовлетворяют уравнению


\begin{equation}


(\chi^2_1\chi^n_2)'+\varepsilon b


\biggl(\frac{\varepsilon b}{2(n+1)}-b^{-1}\biggr) J\chi'_1-


\varepsilon l\chi_1-\mu\chi^{n+1}_2=0,\quad


\chi_2=\chi_1\chi'_1.


\end{equation}


В случае $\mu=0$, $\varepsilon>0$,


$b=\pm(2nl/\varepsilon)^{1/2}$ уравнение (8.8) имеет замкнутое


решение, в


результате для этих значений $\mu$ и $b$ получаем замкнутое решение


уравнения (8.7)


\begin{gather}


u=


\begin{cases}


c_n(1+\varepsilon bt^{1/2})^{-1/n}\exp(2lt^{1/2}/b)


(J^p-\xi^p_0)^{l-s}, & |J|\ge \xi_0;\\


\quad 0, & |J|<\xi_0,


\end{cases}


\\


c_n=(\varepsilon l)^{1/n}l^{-(l-s)}.\notag


\end{gather}


Рассмотрим два случая.





1) $b>0$. Из (8.9) легко следует, что пленка распадается на две


части:


правую


$$


x\ge x^п_f=[\xi_0(1+\varepsilon bt^{1/2})^{-1/n}+


2dnb^{-1}]\exp(2t^{1/2}b^{-1})


$$


и левую


$$


x\le x^л_f=[-\xi_0(1+\varepsilon bt^{1/2})^{-1/n}+


2dnb^{-1}]\exp(2t^{1/2}b^{-1}).


$$


На больших временах ($t\to\infty$) в случае $d>0\;$


$x^{п,л}_f\to\infty$,


т.е. правая половина пленки


отступает, а левая наступает; в случае $d<0$ имеем обратную картину.


Этот


результат нетрудно объяснить, рассматривая второе слагаемое (основное


при больших $t$) в балансе массы в уравнении (8.7). Так как


$\sigma u^{-1}=u_x$ больше


нуля при $x>x^п_f$ и меньше нуля при $x<x^л_f$, то второе слагаемое в


балансе при $\varepsilon>0$,


$d>0$ на решении (8.9) играет роль источника на левой половине


поверхности


пленки и стока на правой; при $\varepsilon>0$, $d<0$ --- наоборот,


имеем источник справа и сток слева.





2) $b<0$. Обозначив $t_0=(\varepsilon b)^{-2}$, найдем, что при


$t\to t_0-0$\; $x^п_f\to\infty$, $x^л_f\to -\infty$,


т.е. за конечное время


правая половина пленки с быстро возрастающей скоростью уходит в


$\infty$, а


левая --- в $-\infty$. Из (8.9) следует, что пленка с ростом $t$


пучится. Этот


результат объясняется тем, что в балансе массы основную роль играет


первое слагаемое


$$


-\varepsilon bt^{-1/2}x\sigma(2(n+1)u)^{-1},


$$


которое для обеих половин поверхности пленки играет при $b<0$,


$\varepsilon>0$ роль поверхностного источника.





Пользуясь случаем, авторы искренне благодарят рецензентов статьи [1] и


данной работы за стимулирующие замечания и А.Н.\,Саламатина и


В.А.\,Чугунова за внимание и обсуждения полученных


результатов.





\begin{thebibliography}{9}





\bibitem{T1}


Тонконог С.Л., Эскин Л.Д. {\em


О некоторых симметриях и инвариантных решениях уравнения


динамики неньютоновской жидкости. I} // Изв. вузов. Математика. --


1998. -- \No\,1. -- С.\,55--66.








\end{thebibliography}





\vskip 2cm





\postup{\it Казанский государственный}{$\qquad$\it Поступили}


\postup{\it университет}{{\it первый вариант} 16.05.1995}


\postup{}{{\it окончательный вариант} 17.04.1997}


\label{end}


\end{document}





